\chapter{System Design and Architecture}

\section{System Architecture}

NutriFit follows a modern three-tier architecture:

\begin{itemize}
    \item \textbf{Presentation Layer}: React-based frontend with Next.js
    \item \textbf{Application Layer}: Next.js API routes handling business logic
    \item \textbf{Data Layer}: SQLite database with Prisma ORM
\end{itemize}

\begin{figure}[H]
\centering
\includegraphics[width=0.7\textwidth]{architecture-diagram.png}
\caption{System Architecture Diagram}
\label{fig:architecture}
\end{figure}

\section{Class Diagram}

\begin{figure}[H]
\centering
\includegraphics[width=0.95\textwidth]{class-diagram.png}
\caption{Core Class Diagram}
\label{fig:class-diagram}
\end{figure}

The class diagram shows the main entities and their relationships:
\begin{itemize}
    \item \textbf{Entity Classes}: User, HealthProfile, Goal, DietPlan, WorkoutPlan, MealLog, ExerciseLog
    \item \textbf{Service Classes}: PlanGenerator, SafetyChecker, ProgressTracker
    \item \textbf{Catalog Classes}: FoodItem, ExerciseItem
    \item \textbf{Relationships}: One-to-one, one-to-many, composition, and dependency relationships
\end{itemize}

\section{Database Design}

\subsection{Entity-Relationship Diagram}

\begin{figure}[H]
\centering
\includegraphics[width=0.9\textwidth]{er-diagram.png}
\caption{Entity-Relationship Diagram}
\label{fig:er-diagram}
\end{figure}

The database consists of 11 main entities with the following relationships:

\begin{itemize}
    \item \textbf{User} (1) --- (0..1) \textbf{HealthProfile}
    \item \textbf{User} (1) --- (*) \textbf{MealLog}
    \item \textbf{User} (1) --- (*) \textbf{ExerciseLog}
    \item \textbf{User} (1) --- (*) \textbf{DietPlan}
    \item \textbf{User} (1) --- (*) \textbf{WorkoutPlan}
    \item \textbf{User} (1) --- (*) \textbf{Goal}
    \item \textbf{User} (1) --- (*) \textbf{WaterIntake}
    \item \textbf{User} (1) --- (*) \textbf{Favorite}
\end{itemize}

\subsection{Key Database Tables}

\subsubsection{User Table}

\begin{table}[H]
\centering
\begin{tabular}{|l|l|l|}
\hline
\textbf{Column} & \textbf{Type} & \textbf{Constraints} \\
\hline
user\_id & INTEGER & PRIMARY KEY \\
name & TEXT & NOT NULL \\
email & TEXT & UNIQUE, NOT NULL \\
password\_hash & TEXT & NOT NULL \\
created\_at & DATETIME & DEFAULT NOW \\
\hline
\end{tabular}
\caption{User Table Schema}
\end{table}

\subsubsection{HealthProfile Table}

\begin{table}[H]
\centering
\begin{tabular}{|l|l|l|}
\hline
\textbf{Column} & \textbf{Type} & \textbf{Constraints} \\
\hline
profile\_id & INTEGER & PRIMARY KEY \\
user\_id & INTEGER & FOREIGN KEY, UNIQUE \\
age & INTEGER & NULL \\
weight & REAL & NULL \\
height & REAL & NULL \\
gender & TEXT & DEFAULT 'male' \\
activity\_level & TEXT & DEFAULT 'moderate' \\
health\_goals & TEXT & DEFAULT 'maintain' \\
diseases & TEXT & JSON array \\
allergens & TEXT & JSON array \\
\hline
\end{tabular}
\caption{HealthProfile Table Schema}
\end{table}

\subsection{Database Indexes}

\begin{table}[H]
\centering
\begin{tabular}{|l|l|p{6cm}|}
\hline
\textbf{Table} & \textbf{Index} & \textbf{Purpose} \\
\hline
MealLog & (userId, date) & Fast date-based queries \\
ExerciseLog & (userId, date) & Fast date-based queries \\
Goal & (userId, status) & Filter active goals \\
WaterIntake & (userId, date) & Daily water tracking \\
Favorite & (userId, itemType) & Filter favorites \\
\hline
\end{tabular}
\caption{Database Indexes}
\end{table}

\section{Technology Stack}

\subsection{Frontend Technologies}

\begin{table}[H]
\centering
\begin{tabular}{|l|l|p{6cm}|}
\hline
\textbf{Technology} & \textbf{Version} & \textbf{Purpose} \\
\hline
Next.js & 14.2.33 & React framework with SSR \\
React & 18.3.1 & UI library \\
TypeScript & 5.4.5 & Type-safe JavaScript \\
Tailwind CSS & 3.4.3 & Utility-first CSS \\
Recharts & 3.2.1 & Chart library \\
\hline
\end{tabular}
\caption{Frontend Stack}
\end{table}

\subsection{Backend Technologies}

\begin{table}[H]
\centering
\begin{tabular}{|l|l|p{6cm}|}
\hline
\textbf{Technology} & \textbf{Version} & \textbf{Purpose} \\
\hline
Node.js & 20.x & JavaScript runtime \\
Prisma & 6.17.0 & ORM \\
bcryptjs & 2.4.3 & Password hashing \\
jose & 5.2.3 & JWT handling \\
SQLite & 3.x & Database \\
\hline
\end{tabular}
\caption{Backend Stack}
\end{table}

\section{Security Design}

\subsection{Authentication Flow}

\begin{enumerate}
    \item User submits credentials
    \item Server validates and hashes password (bcrypt, 10 rounds)
    \item Server generates JWT token (7-day expiration)
    \item Client stores token in localStorage
    \item Client includes token in Authorization header
    \item Server validates token on each request
\end{enumerate}

\subsection{Security Measures}

\begin{itemize}
    \item \textbf{Password Security}: Bcrypt hashing with salt
    \item \textbf{Token Security}: JWT with 7-day expiration
    \item \textbf{API Security}: Token validation on all protected endpoints
    \item \textbf{Input Validation}: Server-side validation for all inputs
    \item \textbf{Error Handling}: Generic error messages (no sensitive data exposure)
\end{itemize}

\section{UI/UX Design}

\subsection{Design Principles}

\begin{itemize}
    \item \textbf{Simplicity}: Clean, uncluttered interface
    \item \textbf{Consistency}: Uniform design patterns
    \item \textbf{Responsiveness}: Mobile-first approach
    \item \textbf{Accessibility}: Clear labels, good contrast
    \item \textbf{Feedback}: Visual confirmation of actions
\end{itemize}

\subsection{Color Scheme}

\begin{itemize}
    \item \textbf{Primary Blue} (\#2563eb): Trust, professionalism
    \item \textbf{Success Green} (\#10b981): Positive actions
    \item \textbf{Warning Yellow} (\#f59e0b): Caution
    \item \textbf{Danger Red} (\#ef4444): Errors, unsafe items
    \item \textbf{Neutral Gray}: Text, backgrounds
\end{itemize}

\subsection{Responsive Breakpoints}

\begin{table}[H]
\centering
\begin{tabular}{|l|l|l|}
\hline
\textbf{Device} & \textbf{Breakpoint} & \textbf{Layout} \\
\hline
Mobile & < 768px & Single column \\
Tablet & 768px - 1024px & 2 columns \\
Desktop & > 1024px & 3-4 columns \\
\hline
\end{tabular}
\caption{Responsive Breakpoints}
\end{table}
